%==============================================================================
%  ENGR 2405 - Electrical Circuits I
%  Lab Report Template
%
%  Compile with pdfLaTeX (or XeLaTeX / LuaLaTeX).
%  If you use the listings code blocks, no extra setup is needed.
%==============================================================================
\documentclass[11pt]{article}

%------------------------- Packages -------------------------------------------
\usepackage[utf8]{inputenc}
\usepackage[T1]{fontenc}
\usepackage[margin=1in]{geometry}     % page margins
\usepackage{amsmath,amssymb}          % math
\usepackage{graphicx}                 % \includegraphics for schematics/scope shots
\usepackage{float}                    % [H] figure placement
\usepackage{booktabs}                 % nicer tables
\usepackage{siunitx}                  % \SI{5}{\ohm}, \SI{10}{\milli\second}, etc.
\usepackage{listings}                 % SciLab / netlist code blocks
\usepackage{xcolor}
\usepackage{caption}
\usepackage[hidelinks]{hyperref}
\usepackage{titlesec}

%------------------------- Section styling ------------------------------------
\titleformat{\section}{\Large\bfseries}{}{0pt}{}
\titlespacing*{\section}{0pt}{1.4em}{0.6em}

%------------------------- Code listing style ---------------------------------
\definecolor{codebg}{RGB}{245,245,245}
\definecolor{codecomment}{RGB}{0,128,0}
\definecolor{codekw}{RGB}{0,0,180}
\lstset{
    backgroundcolor=\color{codebg},
    basicstyle=\ttfamily\small,
    breaklines=true,
    frame=single,
    framesep=4pt,
    keywordstyle=\color{codekw}\bfseries,
    commentstyle=\color{codecomment}\itshape,
    showstringspaces=false,
    captionpos=b,
    tabsize=2,
}
\definecolor{codestr}{RGB}{163,21,21}      % strings
\definecolor{codenum}{RGB}{120,120,120}    % numbers

% --- MATLAB (uses the listings built-in language) -------------------------
\lstdefinestyle{matlab}{
    language=Matlab,
    stringstyle=\color{codestr},
}

% --- SciLab (custom: listings has no built-in SciLab language) -------------
\lstdefinelanguage{Scilab}{
    morekeywords={
        function,endfunction,if,then,else,elseif,end,for,while,do,
        select,case,break,continue,return,clc,clear,disp,mprintf,
        printf,zeros,ones,eye,linspace,inv,det,sum,length,size,abs,
        sqrt,real,imag,exp,sin,cos,tan,plot,plot2d,xlabel,ylabel,
        title,legend,grid,poly,roots,find,pause,error,warning
    },
    sensitive=true,
    morecomment=[l]{//},          % line comment
    morecomment=[s]{/*}{*/},      % block comment
    morestring=[b]",              % double-quoted strings
    morestring=[b]',              % single-quoted strings
}
\lstdefinestyle{scilab}{
    language=Scilab,
    stringstyle=\color{codestr},
}

% --- LTSpice netlist (custom: SPICE directives & comments) -----------------
\lstdefinelanguage{LTspice}{
    morekeywords={
        .tran,.op,.dc,.ac,.step,.param,.model,.lib,.include,.inc,
        .meas,.measure,.save,.backanno,.end,.ends,.subckt,.global,
        .ic,.nodeset,.options,.option,.func,.temp,.four,.noise,.tf
    },
    sensitive=false,
    morecomment=[l]{*},           % lines beginning with * are comments
    morecomment=[l]{;},           % inline ; comments
    morestring=[b]",
}
\lstdefinestyle{netlist}{
    language=LTspice,
    keywordstyle=\color{codekw}\bfseries,
    basicstyle=\ttfamily\small,
}

%==============================================================================
\begin{document}

%------------------------- Title block ----------------------------------------
\begin{center}
    {\LARGE\bfseries SciLab (Lab 02)}\\[0.8em]
    {\large Abdon Morales}\\[0.3em]
    {\large \today}
\end{center}
\vspace{1em}

%------------------------------------------------------------------------------
\section*{Description}
% Enter a brief description of the lab. This can be taken from the lab
% assignment itself.
This lab will familiarize the student with the capabilities and basic use of SciLab which is a
freely available calculation program. It is similar to the much more expensive commercial tool
MATLAB. MATLAB is widely used in industry and has advanced capabilities in graphics
plotting and code generation that go well beyond the capabilities of SciLab.


This psuedo code just describes the math transforms necessary to implement
\begin{enumerate}
	\item Polar to Cartesian transforms and
	\item Cartesian to polar transforms
\end{enumerate}
%------------------------------------------------------------------------------
\section*{Calculations}
% Put in any calculations you made. If you used SciLab, paste your program and
% its runtime printout here. Include any reconciliation equations (e.g. mesh
% currents to branch currents) if they are not part of your SciLab output.

% --- SciLab program ---
\begin{lstlisting}[style=scilab, caption={SciLab Part I}]
%j=%i
function x=polar2cart(m, a); 
x=m*cos(a)+(%j*m*sin(a)); 
endfunction

function p=cart2polar(x);
re = real(x); im = imag(x);
m = sqrt(re^2 + im^2);
a = atan(im, re);
p = [m,a];
endfunction

// ----------------- MAIN PROGRAM -----------------
disp("Abdon Morales");
disp("ENGR 2405 - Polar/Cartesian Conversions Testcases");
disp("========================================");
// --- Test cart2polar ---
disp("--- Testing cart2polar ---");
x1 = 3 + 4*%i;
printf("Input  (cartesian): %g + %gi\n", real(x1), imag(x1));
out1 = cart2polar(x1);
printf("Output [m a]:  m = %g,  a = %g rad  (%g deg)\n", ...
       out1(1), out1(2), out1(2)*180/%pi);
// --- Test polar2cart ---
disp("--- Testing polar2cart ---");
m2 = 5;
a2 = 0.9273;   // ~53.13 deg expressed in radians
printf("Input  (polar): m = %g,  a = %g rad\n", m2, a2);
out2 = polar2cart(m2, a2);
printf("Output (cartesian): %g + %gi\n", real(out2), imag(out2));
// --- Round-trip check  ---
disp("--- Round-trip: cart -> polar -> cart ---");
back = polar2cart(out1(1), out1(2));
printf("Should match 3 + 4i:  %g + %gi\n", real(back), imag(back));
// --- Quadrant test: -3 + 4i (second quadrant) ---
disp("--- Testing cart2polar in Quadrant II ---");
x3 = -3 + 4*%i;
printf("Input  (cartesian): %g + %gi\n", real(x3), imag(x3));
out3 = cart2polar(x3);
printf("Output [m a]:  m = %g,  a = %g rad  (%g deg)\n", ...
       out3(1), out3(2), out3(2)*180/%pi);
\end{lstlisting}

% --- SciLab runtime printout ---
\begin{lstlisting}[caption={SciLab Part I runtime output}]
  "Abdon Morales"
  "ENGR 2405 - Polar/Cartesian Conversions Testcases"
  "========================================"
  "--- Testing cart2polar ---"
Input  (cartesian): 3 + 4i
Output [m a]:  m = 5,  a = 0.927295 rad  (53.1301 deg)
  "--- Testing polar2cart ---"
Input  (polar): m = 5,  a = 0.9273 rad
Output (cartesian): 2.99998 + 4.00001i
  "--- Round-trip: cart -> polar -> cart ---"
Should match 3 + 4i:  3 + 4i
  "--- Testing cart2polar in Quadrant II ---"
Input  (cartesian): -3 + 4i
Output [m a]:  m = 5,  a = 2.2143 rad  (126.87 deg)
\end{lstlisting}

\begin{lstlisting}[style=scilab, caption={SciLab Part II}]
// ----------------- MAIN PROGRAM (Part II) -----------------
disp("Abdon Morales");
disp("ENGR 2405 - System of Linear Equations");
disp("========================================");
// Build the matrix and vector
A = [2 1 1; 1 2 1; 1 1 2];
X = [1;2;3];
disp("Matrix A = "); disp(A);
disp("Vector X = "); disp(X);

// Do the math via matrix-vector multiplication
Y = A*X;
disp("This is the result of A*X");
disp(Y);

// Do the matrix A inversion
invA = inv(A);
disp("This is A inverse =");
disp(invA);

// Now go back to find X using the inversion
out = invA * Y;
disp("Vector X using A inverse * Y");
disp(out);
\end{lstlisting}

\begin{lstlisting}[caption={SciLab Part II runtime output}]
  "Abdon Morales"
  "ENGR 2405 - System of Linear Equations"
  "========================================"
  "Matrix A = "
   2.   1.   1.
   1.   2.   1.
   1.   1.   2.
  "Vector X = "
   1.
   2.
   3.
  "This is the result of A*X"
   7.
   8.
   9.
  "This is A inverse ="
   0.75  -0.25  -0.25
  -0.25   0.75  -0.25
  -0.25  -0.25   0.75
  "Vector X using A inverse * Y"
   1.
   2.
   3.

\end{lstlisting}

%------------------------------------------------------------------------------
\section*{Observations \& Discussion}
% List any anomalies or interesting observations. Discuss any discrepancies and
% their reasons. Answer any questions from the lab.

None.

%==============================================================================
\end{document}
